\chapter{Selected Derivations}
\label{app:derivations}

\section{Normalization of the Maxwell--J\"uttner distribution}

For an isotropic distribution,
\begin{equation}
 n_e=4\pi\int_0^\infty p^2 f(p)\,\dd p.
\end{equation}
Let $u=p/(m_ec)$ and $\gamma=\sqrt{1+u^2}$. The integral identity
\begin{equation}
\int_0^\infty u^2e^{-\gamma/\theta}\,\dd u
=\theta K_2(1/\theta)
\end{equation}
leads to Equation~\eqref{eq:mj}. The energy density is
\begin{equation}
U_e=4\pi\int_0^\infty p^2(\gamma-1)m_ec^2f(p)\,\dd p.
\end{equation}

\section{Number-conserving redistribution}

Starting from normalized $f_0(E)$, define
\begin{equation}
N_>=\int_{E_c}^{\infty}f_0(E)\dd E.
\end{equation}
For a normalized redistribution kernel $g_r$, the new distribution is
\begin{equation}
f_c(E)=f_0(E)H(E_c-E)+N_>g_r(E).
\end{equation}
Then
\begin{align}
\int_0^\infty f_c(E)\dd E
&=\int_0^{E_c}f_0(E)\dd E+N_>\\
&=1-N_>+N_>=1.
\end{align}
The energy change is
\begin{equation}
\Delta\langle E\rangle
=N_>\langle E\rangle_r-\int_{E_c}^{\infty}Ef_0(E)\dd E,
\end{equation}
which is negative when the redistribution mean lies below the original tail mean.

\section{Dose ratio as a weighted spectral average}

Let the controlled source be $S_c(E)=R_S(E)S_0(E)$. For a fixed linear transport and dose weighting $W(E)$,
\begin{equation}
D_0=\int W(E)S_0(E)\dd E,
\qquad
D_c=\int W(E)R_S(E)S_0(E)\dd E.
\end{equation}
Therefore
\begin{equation}
\frac{D_c}{D_0}
=\int R_S(E)\pi_D(E)\dd E,
\end{equation}
where
\begin{equation}
\pi_D(E)=\frac{W(E)S_0(E)}{\int W(E')S_0(E')\dd E'}
\end{equation}
is a normalized nonnegative weighting density. Thus $D_c/D_0$ lies between the minimum and maximum of $R_S$ over the contributing energy range. It equals the total-power ratio only for special weightings or spectrally uniform suppression.

\section{Linear-quadratic survival under dose reduction}

For baseline and controlled doses $D_0$ and $D_c$,
\begin{equation}
\ln\left(\frac{S_c}{S_0}\right)
=-\alpha D_c-\beta D_c^2+\alpha D_0+\beta D_0^2.
\end{equation}
Factorization gives
\begin{equation}
\ln\left(\frac{S_c}{S_0}\right)
=(D_0-D_c)\left[\alpha+\beta(D_0+D_c)\right].
\end{equation}
For nonnegative $\alpha$, $\beta$, and $D_0>D_c$, the controlled exposure yields greater survival. The incremental benefit increases with baseline dose in the quadratic regime.

\section{A simple cutoff criterion}

Suppose collisional replenishment of a tail state occurs on time $\tau_c(E)$ and perturbation loss on $\tau_L(E)$. A local stationary balance has the schematic form
\begin{equation}
0\approx\frac{f_{\mathrm{eq}}-f}{\tau_c}-\frac{f}{\tau_L}.
\end{equation}
Solving,
\begin{equation}
\frac{f}{f_{\mathrm{eq}}}
=\frac{1}{1+\tau_c/\tau_L}.
\end{equation}
The population is reduced by one half where $\tau_c=\tau_L$. This provides a convenient operational definition of the cutoff onset, while a full kinetic model determines the smooth shape.

\section{Mirror loss boundary with electrostatic potential}

Conservation of total energy and magnetic moment gives, between a field minimum and mirror throat,
\begin{equation}
\gamma_0m_ec^2-e\phi_0
=\gamma_mm_ec^2-e\phi_m,
\end{equation}
with $p_{\perp,m}^2/B_m=p_{\perp,0}^2/B_0$ under adiabatic motion. Reflection occurs when $p_{\parallel,m}=0$. Combining these relations yields the relativistic loss boundary in $(p,\xi)$; the nonrelativistic limit reduces to Equation~\eqref{eq:mirror-potential}. Numerical orbit following is preferred when fields vary rapidly or adiabatic invariance fails.
